\section*{Nền tảng cơ học lượng tử}
\begin{longtable}{@{}>{\raggedright\arraybackslash}p{0.21\textwidth}>{\raggedright\arraybackslash\itshape}p{0.21\textwidth}>{\raggedright\arraybackslash}p{0.52\textwidth}@{}}
\toprule \textbf{Tiếng Việt} & \textbf{\upshape English} & \textbf{Giải thích} \\ \midrule \endhead
Hàm sóng & Wavefunction ($\Psi$) & Hàm chứa toàn bộ thông tin về trạng thái lượng tử của hệ. Với N electron, $\Psi$ phụ thuộc đồng thời vào tọa độ và spin của cả N electron (Li có 3 electron, tức 9 tọa độ không gian). \\[2pt]
Phương trình Schrödinger dừng & Time-independent Schrödinger equation & Phương trình trị riêng $\hat H\Psi$ = E$\Psi$. Giải nó cho ra các mức năng lượng E được phép và hàm sóng tương ứng; với từ 2 electron trở lên thì không có nghiệm giải tích. \\[2pt]
Toán tử Hamilton (Hamiltonian) & Hamiltonian operator ($\hat H$) & Toán tử ứng với năng lượng toàn phần, tức động năng cộng thế năng. Với nguyên tử, nó gồm động năng các electron, lực hút hạt nhân–electron và lực đẩy giữa các electron. \\[2pt]
Trị riêng; hàm riêng / vectơ riêng & Eigenvalue; eigenfunction / eigenvector & Nếu $\hat A$f = a$\cdot$f thì f là hàm riêng và a là trị riêng của $\hat A$. Khi khai triển theo một tập cơ sở, bài toán trở thành chéo hóa ma trận trên máy tính. \\[2pt]
Toán tử Hermite (tự liên hợp) & Hermitian (self-adjoint) operator & Toán tử có trị riêng luôn là số thực và các hàm riêng trực giao với nhau. Mọi đại lượng đo được (năng lượng, động lượng, spin…) đều ứng với một toán tử Hermite. \\[2pt]
Mật độ xác suất; chuẩn hóa & Probability density; normalization & |$\Psi$|$^{2}$ là mật độ xác suất tìm thấy hệ ở một cấu hình nhất định (diễn giải Born). Chuẩn hóa nghĩa là $\int$|$\Psi$|$^{2}$ = 1, tức tổng xác suất bằng 1. \\[2pt]
Giá trị kỳ vọng (giá trị trung bình) & Expectation value $\langle\hat A\rangle$ & Giá trị trung bình thu được khi đo đại lượng A rất nhiều lần trên các hệ giống hệt nhau, $\langle\hat A\rangle$ = $\langle\Psi$|$\hat A$|$\Psi\rangle$. $\langle\hat H\rangle$ là năng lượng trung bình, đại lượng mà mọi phương pháp biến phân tìm cách cực tiểu hóa. \\[2pt]
Không gian Hilbert; nguyên lý chồng chất; ký hiệu Dirac & Hilbert space; superposition principle; Dirac (bra–ket) notation & Trạng thái lượng tử là các ``vectơ'' trong một không gian có tích vô hướng, và tổ hợp tuyến tính của các trạng thái cũng là một trạng thái hợp lệ. Ký hiệu |$\psi\rangle$, $\langle\varphi$|$\psi\rangle$ giúp viết gọn; mọi tài liệu từ G5 trở đi đều dùng. \\[2pt]
Trạng thái cơ bản / trạng thái kích thích & Ground state / excited state & Trạng thái có năng lượng thấp nhất / các trạng thái có năng lượng cao hơn. Báo cáo này chủ yếu tìm năng lượng trạng thái cơ bản của Li (cấu hình 1s$^{2}$2s, ký hiệu $^{2}$S). \\[2pt]
Đơn vị nguyên tử & Atomic units (hartree $E_h$, bohr $a_0$) & Hệ đơn vị đặt $\hbar$ = $m_e$ = e = 4$\pi\varepsilon_0$ = 1 để phương trình gọn hơn. 1 $E_h$ $\approx$ 27.211 eV $\approx$ 627.5 kcal/mol; 1 $a_0$ $\approx$ 0.529 Å. \\[2pt]
Spin & Spin & Mômen động lượng nội tại của electron (s = 1/2), có hai trạng thái `\texttt{lên'' ($\alpha$) và }`xuống'' ($\beta$). Spin không có mặt trong Hamiltonian phi tương đối tính nhưng chi phối cách xếp electron thông qua nguyên lý Pauli. \\[2pt]
Orbital (orbital); spin-orbital & Orbital; spin orbital & Orbital là hàm sóng của một electron (1s, 2s, 2p…); spin-orbital là orbital nhân với hàm spin $\alpha$ hoặc $\beta$. Spin-orbital là ``viên gạch'' để xây định thức Slater. \\[2pt]
Cấu hình electron & Electron configuration & Cách phân bố electron vào các orbital, ví dụ Li là 1s$^{2}$2s$^{1}$. Đây chỉ là bức tranh gần đúng: hàm sóng thật là tổ hợp của nhiều cấu hình. \\[2pt]
Nguyên lý biến phân & Variational principle & Với mọi hàm thử hợp lệ, $\langle\hat H\rangle$ $\ge$ E$_{0}$ (năng lượng cơ bản đúng). Vì vậy ``năng lượng thấp hơn nghĩa là hàm thử tốt hơn''; đây là nền tảng của HF, CI, VMC, FermiNet và VQE. \\[2pt]
Hàm thử; dạng thử (ansatz) & Trial wavefunction; ansatz & Một dạng hàm sóng giả định có chứa tham số điều chỉnh được, ví dụ $e^{-\zeta(r_1+r_2)}$, một định thức Slater, một mạng nơ-ron hay một mạch lượng tử. Ansatz càng linh hoạt thì kết quả càng sát nhưng càng tốn chi phí tính. \\[2pt]
Hiệu ứng chắn; điện tích hạt nhân hiệu dụng & Screening; effective nuclear charge ($Z_{\mathrm{eff}}$) & Các electron bên trong ``che'' bớt hạt nhân, nên electron bên ngoài chỉ cảm thấy một điện tích nhỏ hơn Z. Ví dụ electron 2s của Li cảm thấy điện tích nhỏ hơn 3 rất nhiều. \\[2pt]
Năng lượng ion hóa & Ionization energy (IE) & Năng lượng tối thiểu để tách một electron khỏi nguyên tử. Với Li là 5.3917 eV (NIST); đây là phép thử nhạy cho mọi phương pháp vì nó là hiệu của hai năng lượng lớn. \\[2pt]
Ký hiệu số hạng & Term symbol ($^{2S+1}L_J$) & Ký hiệu gọn cho spin toàn phần S, mômen động lượng quỹ đạo L và mômen toàn phần J của một trạng thái nguyên tử. Trạng thái cơ bản của Li là $^2S_{1/2}$. \\[2pt]
\bottomrule\end{longtable}
\section*{Bài toán nhiều electron và tính phản đối xứng}
\begin{longtable}{@{}>{\raggedright\arraybackslash}p{0.21\textwidth}>{\raggedright\arraybackslash\itshape}p{0.21\textwidth}>{\raggedright\arraybackslash}p{0.52\textwidth}@{}}
\toprule \textbf{Tiếng Việt} & \textbf{\upshape English} & \textbf{Giải thích} \\ \midrule \endhead
Bài toán nhiều hạt (nhiều vật) & Many-body problem & Bài toán có nhiều hạt tương tác với nhau, nên chuyển động của mỗi hạt phụ thuộc vào mọi hạt khác. Chỉ với 3 electron (Li) đã không còn nghiệm chính xác dạng giải tích. \\[2pt]
Gần đúng Born–Oppenheimer & Born–Oppenheimer approximation & Hạt nhân nặng hơn electron hàng nghìn lần nên được coi là đứng yên, và ta chỉ giải phương trình cho electron. Với một nguyên tử, điều này tương đương coi hạt nhân có khối lượng vô hạn. \\[2pt]
Hamiltonian electron; lực đẩy electron–electron & Electronic Hamiltonian; electron–electron repulsion & Trong đơn vị nguyên tử, $\hat H$ = $\Sigma_{i}$(-½$\nabla_{i}^{2}$ - Z/r$_{i}$) + $\Sigma$\_\{i<j\} 1/r$_{ij}$. Số hạng 1/r$_{ij}$ ràng buộc các electron với nhau, khiến phương trình không tách được thành các bài toán một electron. \\[2pt]
Hạt đồng nhất; fermion và boson & Identical particles; fermions and bosons & Các electron giống hệt nhau nên không thể ``đánh số'' để theo dõi từng hạt. Hạt có spin bán nguyên (như electron) là fermion, hàm sóng đổi dấu khi hoán đổi hai hạt; hạt có spin nguyên là boson, hàm sóng không đổi dấu. \\[2pt]
Tính phản đối xứng & Antisymmetry & Hàm sóng electron phải đổi dấu khi hoán đổi tọa độ (cả không gian lẫn spin) của hai electron bất kỳ. Đây là ràng buộc then chốt mà mọi ansatz, từ định thức Slater đến FermiNet, đều phải thỏa mãn. \\[2pt]
Nguyên lý loại trừ Pauli & Pauli exclusion principle & Không thể có hai electron cùng nằm trong một spin-orbital; đây là hệ quả trực tiếp của tính phản đối xứng. Vì vậy electron thứ ba của Li phải lên 2s, cấu hình 1s$^{3}$ bị cấm. \\[2pt]
Tích Hartree; định thức Slater & Hartree product; Slater determinant & Tích Hartree $\varphi_{1}$(x$_{1}$)$\varphi_{2}$(x$_{2}$)$\varphi_{3}$(x$_{3}$) đơn giản nhưng không phản đối xứng. Định thức Slater (định thức của các spin-orbital, chia cho $\sqrt{}$N!) tự động phản đối xứng và bằng 0 khi hai electron rơi vào cùng một spin-orbital. \\[2pt]
Lỗ trao đổi (lỗ Fermi); lỗ Coulomb & Exchange (Fermi) hole; Coulomb hole & Vùng quanh mỗi electron nơi xác suất gặp electron khác bị giảm: do nguyên lý Pauli (với electron cùng spin, HF đã mô tả được) hoặc do lực đẩy Coulomb (với mọi spin, HF bỏ sót). Phiếm hàm DFT thực chất là mô hình hóa các ``lỗ'' này. \\[2pt]
Điều kiện điểm nhọn Kato (cusp) & Kato cusp condition & Hàm sóng đúng có ``điểm gãy'' (đạo hàm có bước nhảy xác định) khi một electron chạm hạt nhân hoặc hai electron chạm nhau. Các hàm Gauss trơn không tái tạo được điểm gãy này nên cần rất nhiều hàm cơ sở; trong QMC, thừa số Jastrow xử lý trực tiếp. \\[2pt]
Lời nguyền số chiều & Curse of dimensionality & Hàm sóng N electron sống trong không gian 3N chiều. Lưu nó trên lưới 100 điểm mỗi chiều cần 100\textasciicircum{}(3N) số, với Li đã là 10$^{18}$ số (khoảng 8 exabyte). Mọi phương pháp trong báo cáo đều là cách ``né'' lời nguyền này. \\[2pt]
Ma trận mật độ rút gọn & Reduced density matrix (1-RDM, 2-RDM) & ``Bản nén'' của hàm sóng, chỉ giữ thông tin về 1 hoặc 2 electron; vẫn đủ để tính năng lượng vì Hamiltonian chỉ chứa toán tử một hạt và hai hạt. Mật độ electron n(r) trong DFT chính là phần đường chéo của 1-RDM. \\[2pt]
\bottomrule\end{longtable}
\section*{Lượng tử hóa lần thứ hai}
\begin{longtable}{@{}>{\raggedright\arraybackslash}p{0.21\textwidth}>{\raggedright\arraybackslash\itshape}p{0.21\textwidth}>{\raggedright\arraybackslash}p{0.52\textwidth}@{}}
\toprule \textbf{Tiếng Việt} & \textbf{\upshape English} & \textbf{Giải thích} \\ \midrule \endhead
Lượng tử hóa lần thứ hai & Second quantization & Ngôn ngữ mô tả trạng thái bằng câu hỏi `\texttt{orbital nào đang bị chiếm'' thay vì }`electron nào ở đâu''. Tính phản đối xứng được cài sẵn trong đại số toán tử, nên các công thức gọn hơn rất nhiều. \\[2pt]
Toán tử sinh / toán tử hủy; hệ thức phản giao hoán & Creation / annihilation operators; anticommutation relations & a$^\dagger_{p}$ thêm một electron vào spin-orbital p, còn a$_{p}$ lấy đi một electron. Chúng thỏa \{a$_{p}$, a$^\dagger_{q}$\} = $\delta_{pq}$ và \{a$_{p}$, a$_{q}$\} = 0; từ đó suy ra (a$^\dagger_{p}$)$^{2}$ = 0, chính là nguyên lý Pauli. \\[2pt]
Không gian Fock; vectơ số chiếm & Fock space; occupation-number vector & Không gian chứa các trạng thái với số electron bất kỳ. Một định thức Slater được viết gọn như |1,1,1,0,0,…$\rangle$ (1 là orbital có electron, 0 là orbital trống); đây cũng chính là cách mã hóa trạng thái trên qubit. \\[2pt]
Quy tắc Slater–Condon & Slater–Condon rules & Công thức tính phần tử ma trận $\langle\Phi_{I}$|$\hat H$|$\Phi_{J}\rangle$ giữa hai định thức Slater; nó chỉ khác 0 khi hai định thức khác nhau nhiều nhất hai spin-orbital. Đây là ``động cơ'' của mọi chương trình CI. \\[2pt]
\bottomrule\end{longtable}
\section*{Phương pháp hàm sóng: Hartree–Fock và hậu HF}
\begin{longtable}{@{}>{\raggedright\arraybackslash}p{0.21\textwidth}>{\raggedright\arraybackslash\itshape}p{0.21\textwidth}>{\raggedright\arraybackslash}p{0.52\textwidth}@{}}
\toprule \textbf{Tiếng Việt} & \textbf{\upshape English} & \textbf{Giải thích} \\ \midrule \endhead
Phương pháp Hartree–Fock (HF); gần đúng trường trung bình & Hartree–Fock method; mean-field approximation & Tìm một định thức Slater có năng lượng thấp nhất, trong đó mỗi electron chuyển động trong ``trường trung bình'' của các electron còn lại. HF thu được khoảng 99\% năng lượng tổng nhưng bỏ sót tương quan. \\[2pt]
Trường tự hợp (SCF) & Self-consistent field & Vòng lặp: đoán orbital, tính trường trung bình, giải lại orbital, rồi lặp đến khi kết quả không đổi (tức là ``tự hợp''). Dùng cho cả HF và Kohn–Sham DFT; DIIS là kỹ thuật tăng tốc hội tụ phổ biến. \\[2pt]
Toán tử Fock & Fock operator (f̂) & Hamiltonian hiệu dụng cho một electron trong HF: động năng, cộng lực hút hạt nhân, cộng thế Coulomb trung bình (J), trừ thế trao đổi (K). Orbital HF là hàm riêng của toán tử này, năng lượng orbital là trị riêng tương ứng. \\[2pt]
Tích phân Coulomb (J), tích phân trao đổi (K); năng lượng trao đổi & Coulomb and exchange integrals; exchange energy & J là lực đẩy cổ điển giữa hai ``đám mây'' electron. K là hiệu chỉnh thuần lượng tử, chỉ có giữa các electron cùng spin, và làm năng lượng giảm; năng lượng trao đổi không có tương tự trong vật lý cổ điển. \\[2pt]
Phương trình Roothaan–Hall; ma trận xen phủ & Roothaan–Hall equations (FC = SC$\varepsilon$); overlap matrix S & Dạng ma trận của phương trình HF khi khai triển orbital theo một tập cơ sở. S\_$\mu\nu$ = $\langle\chi$\_$\mu$|$\chi$\_$\nu\rangle$ xuất hiện vì các hàm cơ sở không trực giao với nhau; phương trình được giải bằng chéo hóa ma trận tổng quát. \\[2pt]
HF hạn chế / hạn chế vỏ mở / không hạn chế; nhiễm bẩn spin & RHF / ROHF / UHF; spin contamination & RHF dùng chung một orbital không gian cho cặp $\alpha$–$\beta$ (hệ vỏ kín). Li có một electron độc thân nên cần ROHF hoặc UHF. UHF cho phép orbital $\alpha$ và $\beta$ khác nhau nhưng có thể bị ``nhiễm bẩn spin'', tức $\langle$Ŝ$^{2}\rangle$ lệch khỏi giá trị đúng (0.75 với Li). \\[2pt]
Định lý Koopmans & Koopmans' theorem & Trong HF, năng lượng ion hóa xấp xỉ bằng -$\varepsilon_{HOMO}$ (năng lượng orbital bị chiếm cao nhất, đổi dấu). Đây là ước lượng nhanh nhưng bỏ qua sự ``hồi phục'' của các orbital và hiệu ứng tương quan. \\[2pt]
Giới hạn Hartree–Fock & Hartree–Fock limit & Năng lượng HF khi dùng tập cơ sở đầy đủ, mức thấp nhất mà một định thức Slater có thể đạt tới. Với Li là $\approx$ -7.432727 $E_h$ (Bunge et al. 1993). \\[2pt]
Bộ hàm cơ sở (tập cơ sở) & Basis set & Tập hàm đã biết (thường đặt tại các hạt nhân) dùng để khai triển orbital, biến phương trình vi phân thành bài toán ma trận. Tập càng lớn thì càng chính xác nhưng càng tốn: số tích phân hai electron tăng cỡ M$^{4}$ với M hàm cơ sở. \\[2pt]
Hàm kiểu Slater (STO); hàm kiểu Gauss (GTO); hàm rút gọn & Slater-type orbital; Gaussian-type orbital; contracted Gaussians & STO có dạng e\textasciicircum{}(-$\zeta$r), đúng về mặt vật lý (có điểm gãy cusp). GTO có dạng e\textasciicircum{}(-$\alpha$r$^{2}$), kém hơn về hình dạng nhưng tích phân tính rất nhanh; phần mềm dùng các tổ hợp cố định của nhiều GTO (hàm ``rút gọn'') để bắt chước STO. \\[2pt]
Các họ tập cơ sở; hàm phân cực, hàm khuếch tán; giới hạn tập cơ sở đầy đủ (CBS) & STO-3G, 6-31G\*, cc-pVXZ, cc-pCVXZ, aug-cc-pVXZ; polarization / diffuse functions; CBS limit & STO-3G là tập tối thiểu. cc-pVXZ (X = D, T, Q, 5) tăng dần một cách có hệ thống nên có thể ngoại suy về giới hạn CBS. Chữ `\texttt{C'' nghĩa là có thêm hàm cho tương quan lõi (quan trọng với cặp 1s$^{2}$ của Li); }\texttt{aug-'' thêm hàm khuếch tán; hàm phân cực (d, f…) cho phép orbital }`biến dạng''. \\[2pt]
Năng lượng tương quan & Correlation energy & E$_{c}$ = E$_{\text{đúng}}$ - E$_{HF}$ (định nghĩa của Löwdin), là phần HF bỏ sót vì trong thực tế các electron chủ động né nhau. Với Li, E$_{c}$ $\approx$ -0.0453 $E_h$ $\approx$ -1.23 eV: chỉ khoảng 0.6\% năng lượng tổng nhưng bằng chừng một phần tư năng lượng ion hóa. \\[2pt]
Tương quan động / tương quan tĩnh & Dynamic / static (nondynamic) correlation & Tương quan động là việc electron né nhau ở khoảng cách ngắn, cần rất nhiều kích thích nhỏ để mô tả (MP2 và CC làm tốt). Tương quan tĩnh xảy ra khi vài cấu hình gần suy biến cùng quan trọng (phá liên kết, kim loại chuyển tiếp), cần MCSCF hoặc DMRG. \\[2pt]
Tương tác cấu hình (CI); CI đầy đủ (FCI); định thức kích thích & Configuration interaction; full CI; excited determinants & Hàm sóng được viết là tổ hợp tuyến tính của định thức HF và các định thức ``kích thích'' (chuyển 1, 2, 3… electron sang orbital trống). FCI dùng tất cả các định thức nên chính xác trong tập cơ sở đã chọn, nhưng số định thức tăng theo hàm mũ; CISD chỉ giữ kích thích đơn và đôi. \\[2pt]
Tính nhất quán kích thước; tính quảng tính kích thước & Size consistency; size extensivity & Năng lượng của hai hệ ở rất xa nhau phải bằng tổng năng lượng từng hệ, và năng lượng tương quan phải tăng đúng tỉ lệ với kích thước hệ. MP2 và CC thỏa mãn các yêu cầu này; CISD (CI bị cắt cụt) thì không. \\[2pt]
Lý thuyết nhiễu loạn Møller–Plesset (MP2) & Møller–Plesset perturbation theory & Coi phần chênh lệch giữa Hamiltonian thật và Hamiltonian HF là một ``nhiễu loạn'' nhỏ; MP2 là hiệu chỉnh bậc hai. Đây là cách rẻ nhất (chi phí cỡ N$^{5}$) để thêm tương quan động. \\[2pt]
Phương pháp chùm tương tác (coupled cluster, CC); toán tử cụm & Coupled-cluster method; cluster operator T & Viết $\Psi$ = e\textasciicircum{}T $\Phi_{HF}$, trong đó T = T$_{1}$ + T$_{2}$ + … sinh ra các kích thích đơn, đôi… Nhờ dạng mũ, các tích như T$_{2}^{2}$ xuất hiện tự động, nên CC nhất quán kích thước. \\[2pt]
CCSD(T), ``tiêu chuẩn vàng'' & CCSD(T), ``gold standard'' & CCSD cộng thêm hiệu chỉnh nhiễu loạn cho kích thích bậc ba (T). Rất chính xác cho phân tử vỏ kín gần cấu hình cân bằng, chi phí cỡ N$^{7}$. Với hệ đúng 3 electron như Li, CCSDT đã trùng với FCI. \\[2pt]
MCSCF / CASSCF; không gian hoạt động & Multiconfigurational SCF / complete active space SCF; active space & Tối ưu đồng thời các orbital và hệ số của nhiều cấu hình trong một ``không gian hoạt động'' (một số electron và orbital được chọn trước). Đây là công cụ chính cho tương quan tĩnh; thường được bổ sung CASPT2 hoặc NEVPT2 để thêm tương quan động. \\[2pt]
Hàm sóng tương quan tường minh (Hylleraas, F12) & Explicitly correlated wavefunctions (Hylleraas, F12) & Đưa thẳng khoảng cách r$_{ij}$ giữa các electron vào hàm sóng, nhờ đó mô tả điểm gãy electron–electron tốt hơn nhiều. Phương pháp Hylleraas cho năng lượng Li chính xác đến khoảng 10$^{-12}$ $E_h$ nhưng chỉ khả thi với vài electron; F12 là phiên bản thực dụng cho phân tử. \\[2pt]
Thang chi phí tính toán & Computational scaling & Cách thời gian tính tăng theo kích thước N của hệ: HF cỡ N$^{4}$ (về hình thức), MP2 cỡ N$^{5}$, CCSD cỡ N$^{6}$, CCSD(T) cỡ N$^{7}$, còn FCI tăng theo hàm mũ. Ví dụ, gấp đôi kích thước hệ với CCSD(T) thì thời gian tăng khoảng 2$^{7}$ = 128 lần. \\[2pt]
Độ chính xác hóa học & Chemical accuracy & Sai số không quá 1 kcal/mol ($\approx$ 1.6 m$E_h$ $\approx$ 0.043 eV), mục tiêu thường đặt ra cho dự đoán năng lượng phản ứng. Lưu ý mục tiêu này áp dụng cho năng lượng tương đối (hiệu năng lượng), không phải năng lượng tổng. \\[2pt]
Ab initio (từ nguyên lý đầu) & Ab initio / first principles & Tính toán chỉ dựa trên các hằng số cơ bản và phương trình Schrödinger, không làm khớp tham số với thực nghiệm. HF, CI, CC, QMC đều là ab initio; DFT thường được xếp vào nhóm ``nguyên lý đầu'', nhưng một số phiếm hàm có chứa tham số làm khớp. \\[2pt]
\bottomrule\end{longtable}
\section*{Lý thuyết phiếm hàm mật độ}
\begin{longtable}{@{}>{\raggedright\arraybackslash}p{0.21\textwidth}>{\raggedright\arraybackslash\itshape}p{0.21\textwidth}>{\raggedright\arraybackslash}p{0.52\textwidth}@{}}
\toprule \textbf{Tiếng Việt} & \textbf{\upshape English} & \textbf{Giải thích} \\ \midrule \endhead
Lý thuyết phiếm hàm mật độ (DFT) & Density functional theory & Dùng mật độ electron n(r) (hàm của 3 biến) thay cho hàm sóng (hàm của 3N biến) làm đại lượng cơ bản. Đây là phương pháp tính cấu trúc electron được dùng nhiều nhất hiện nay nhờ cân bằng tốt giữa chi phí và độ chính xác. \\[2pt]
Phiếm hàm & Functional & ``Hàm của hàm'': nhận vào cả một hàm (ví dụ n(r)) và trả ra một con số (ví dụ năng lượng E[n]). Khác với hàm số thông thường, vốn nhận vào một con số. \\[2pt]
Mật độ electron & Electron density n(r) & Số electron trung bình trên một đơn vị thể tích tại điểm r; tích phân trên toàn không gian bằng N (với Li là 3). Khác với hàm sóng, mật độ electron đo được bằng thực nghiệm, ví dụ qua nhiễu xạ tia X. \\[2pt]
Thế ngoài; định lý Hohenberg–Kohn & External potential; Hohenberg–Kohn theorems & Thế ngoài v(r) là thế do hạt nhân tác dụng lên electron (với Li là -3/r). Định lý Hohenberg–Kohn (1964) khẳng định mật độ ở trạng thái cơ bản xác định duy nhất v(r), do đó xác định mọi tính chất của hệ, và tồn tại một phiếm hàm năng lượng mà cực tiểu của nó cho đúng năng lượng cơ bản. Tuy vậy định lý không cho biết phiếm hàm đó có dạng thế nào. \\[2pt]
Phương trình Kohn–Sham; orbital Kohn–Sham & Kohn–Sham equations; Kohn–Sham orbitals & Kohn và Sham (1965) thay hệ electron tương tác bằng một hệ ``giả'' gồm các electron không tương tác nhưng có cùng mật độ, rồi giải tự hợp giống như HF. Orbital Kohn–Sham là công cụ để dựng mật độ và động năng, không phải hàm sóng thật của electron. \\[2pt]
Năng lượng trao đổi–tương quan & Exchange–correlation energy (E$_{xc}$) & Thành phần ``chưa biết'' duy nhất trong Kohn–Sham DFT, gom mọi hiệu ứng lượng tử nhiều hạt (trao đổi, tương quan và phần động năng còn thiếu). Mỗi phiếm hàm gần đúng như LDA, PBE, B3LYP là một cách ước lượng E$_{xc}$. \\[2pt]
Gần đúng mật độ địa phương (LDA) & Local density approximation & Tại mỗi điểm, coi electron như một khí electron đồng nhất có cùng mật độ. Cách này đơn giản và khá tốt cho chất rắn, nhưng làm liên kết trong phân tử mạnh quá mức (overbinding). \\[2pt]
Gần đúng gradient tổng quát (GGA); meta-GGA & Generalized gradient approximation; meta-GGA & GGA dùng thêm gradient của mật độ $\nabla$n (ví dụ PBE, BLYP); meta-GGA dùng thêm mật độ động năng $\tau$ (ví dụ SCAN). Cả hai cải thiện rõ so với LDA trong các bài toán hóa học. \\[2pt]
Phiếm hàm lai (phiếm hàm hỗn hợp) & Hybrid functional & Trộn một phần năng lượng trao đổi ``chính xác'' kiểu HF vào phiếm hàm DFT, ví dụ B3LYP, PBE0. Đây là loại phiếm hàm phổ biến nhất trong hóa học phân tử. \\[2pt]
Thang Jacob & Jacob's ladder & Hình ảnh do J. Perdew đề xuất: xếp các phiếm hàm thành từng bậc LDA $\to$ GGA $\to$ meta-GGA $\to$ phiếm hàm lai $\to$ … theo lượng thông tin được dùng, hướng tới ``thiên đường độ chính xác''. Lên bậc cao thường chính xác hơn, nhưng không được bảo đảm. \\[2pt]
Khí electron đồng nhất; mô hình Thomas–Fermi & Uniform electron gas; Thomas–Fermi model & Khí electron đồng nhất là các electron chuyển động trong một nền điện tích dương phân bố đều; năng lượng của nó (tính rất chính xác bằng QMC) là nền móng của LDA. Mô hình Thomas–Fermi (1927) là mô hình DFT đầu tiên, quá thô (không mô tả được liên kết hóa học) nhưng đã mở đường cho DFT hiện đại. \\[2pt]
Lỗi tự tương tác & Self-interaction error & Trong DFT gần đúng, mỗi electron `\texttt{tự đẩy chính nó'', vì phần Coulomb và phần trao đổi không triệt tiêu nhau hoàn toàn (HF không mắc lỗi này). Hệ quả là -$\varepsilon_{HOMO}$ nhỏ hơn năng lượng ion hóa thực, rào năng lượng phản ứng thấp quá mức, electron bị }`nhòe'' ra quá nhiều. \\[2pt]
Giả thế; thế lõi hiệu dụng (ECP) & Pseudopotential; effective core potential & Thay hạt nhân và các electron lõi bằng một thế hiệu dụng tác dụng lên electron hóa trị, nhờ đó giảm mạnh chi phí tính. Rất phổ biến trong DFT sóng phẳng cho chất rắn và trong QMC. \\[2pt]
\bottomrule\end{longtable}
\section*{Các phương pháp hiện đại: QMC, mạng nơ-ron, máy tính lượng tử, DMRG}
\begin{longtable}{@{}>{\raggedright\arraybackslash}p{0.21\textwidth}>{\raggedright\arraybackslash\itshape}p{0.21\textwidth}>{\raggedright\arraybackslash}p{0.52\textwidth}@{}}
\toprule \textbf{Tiếng Việt} & \textbf{\upshape English} & \textbf{Giải thích} \\ \midrule \endhead
Monte Carlo lượng tử (QMC); sai số thống kê & Quantum Monte Carlo; statistical error & Họ phương pháp dùng số ngẫu nhiên để tính các tích phân nhiều chiều chứa hàm sóng, nên ít bị ảnh hưởng bởi lời nguyền số chiều. Kết quả luôn kèm thanh sai số, giảm như 1/$\sqrt{}$M với M mẫu: muốn sai số nhỏ đi 10 lần thì cần 100 lần số mẫu. \\[2pt]
Monte Carlo biến phân (VMC) & Variational Monte Carlo & Lấy mẫu các cấu hình electron theo phân bố |$\Psi_{T}$|$^{2}$ rồi lấy trung bình năng lượng địa phương; kết quả là một cận trên của năng lượng đúng. VMC là ``động cơ'' bên trong FermiNet và PauliNet. \\[2pt]
Thuật toán Metropolis–Hastings; chuỗi Markov & Metropolis–Hastings algorithm; Markov chain & Cách sinh mẫu theo một phân bố cho trước: đề xuất một bước di chuyển ngẫu nhiên, rồi chấp nhận bước đó với xác suất min(1, P$_{\text{mới}}$/P$_{\text{cũ}}$). Các mẫu liên tiếp tạo thành một chuỗi Markov có tự tương quan, cần tính đến khi ước lượng sai số. \\[2pt]
Năng lượng địa phương & Local energy E$_{L}$ = $\hat H\Psi$/$\Psi$ & Đại lượng tính tại từng cấu hình electron R. Nếu $\Psi$ là hàm riêng đúng thì E$_{L}$ là hằng số và phương sai bằng 0; vì thế phương sai của E$_{L}$ là thước đo chất lượng hàm thử (``nguyên lý phương sai bằng không''). \\[2pt]
Thừa số Jastrow; hàm sóng Slater–Jastrow & Jastrow factor; Slater–Jastrow wavefunction & Nhân định thức Slater với e\textasciicircum{}J, trong đó J phụ thuộc trực tiếp vào khoảng cách electron–electron và electron–hạt nhân, để mô tả tương quan và điều kiện cusp. Hàm Slater–Jastrow là ansatz tiêu chuẩn của QMC. \\[2pt]
Monte Carlo khuếch tán (DMC); thời gian ảo & Diffusion Monte Carlo; imaginary time & Thay t bằng -i$\tau$ biến phương trình Schrödinger thành một phương trình khuếch tán. Mô phỏng bằng các ``walker'' ngẫu nhiên được sinh ra và mất đi; các trạng thái kích thích tắt dần, chỉ còn lại trạng thái cơ bản. DMC chính xác hơn VMC và ít phụ thuộc vào hàm thử hơn. \\[2pt]
Bề mặt nút; gần đúng nút cố định; bài toán dấu & Nodal surface; fixed-node approximation; fermion sign problem & Nút là tập các điểm mà $\Psi$ = 0. Vì hàm sóng fermion đổi dấu, mô phỏng ngẫu nhiên chính xác gặp nhiễu tăng theo hàm mũ (bài toán dấu). DMC ``né'' vấn đề bằng cách giữ cố định nút của hàm thử, nên sai số còn lại chỉ phụ thuộc chất lượng nút: Li có nút do hai electron cùng spin, trong khi trạng thái cơ bản của He thì không có. \\[2pt]
Mạng nơ-ron nhân tạo & Artificial neural network & Một hàm gồm nhiều lớp phép biến đổi tuyến tính xen kẽ với hàm kích hoạt phi tuyến, có rất nhiều tham số; khi đủ lớn, nó xấp xỉ được những hàm phức tạp tùy ý. Mạng được huấn luyện bằng các biến thể của phương pháp hạ gradient. \\[2pt]
Hàm sóng mạng nơ-ron & Neural-network wavefunction & Dùng mạng nơ-ron làm ansatz cho $\Psi$ và tối ưu năng lượng bằng VMC, không cần tập cơ sở Gauss. Với phân tử nhỏ, phương pháp này đạt độ chính xác ngang hoặc tốt hơn CCSD(T), nhưng đòi hỏi GPU. \\[2pt]
FermiNet, PauliNet, Psiformer; đẳng biến hoán vị & (tên riêng); permutation equivariance & FermiNet (DeepMind, 2020) xây các ``định thức'' từ những đặc trưng phụ thuộc mọi electron; PauliNet (2020) cài sẵn kiến thức vật lý (HF, điều kiện cusp); Psiformer (2023) dùng cơ chế tự chú ý. Cả ba được thiết kế đẳng biến hoán vị, để hàm sóng cuối cùng phản đối xứng đúng như yêu cầu. \\[2pt]
Máy tính lượng tử; qubit; NISQ & Quantum computer; qubit; noisy intermediate-scale quantum & Qubit là một hệ hai mức có thể ở trạng thái chồng chất của |0$\rangle$ và |1$\rangle$; máy tính lượng tử thao tác trên qubit bằng các cổng lượng tử. Các máy hiện nay (gọi là NISQ) có số qubit hạn chế, còn nhiễu và chưa được hiệu chỉnh lỗi đầy đủ. \\[2pt]
Rối lượng tử (vướng víu lượng tử) & Quantum entanglement & Tương quan lượng tử khiến không thể mô tả từng phần của hệ một cách riêng rẽ. Rối lượng tử vừa là ``tài nguyên'' của máy tính lượng tử, vừa là thước đo độ khó của bài toán: trong DMRG, entropy rối quyết định số chiều liên kết cần dùng. \\[2pt]
Bộ giải trị riêng lượng tử biến phân (VQE); mạch lượng tử tham số hóa & Variational quantum eigensolver; parameterized quantum circuit & Thuật toán lai: máy lượng tử chuẩn bị trạng thái bằng một mạch có tham số (đóng vai trò ansatz) và đo $\langle\hat H\rangle$; máy tính cổ điển điều chỉnh tham số để cực tiểu năng lượng theo nguyên lý biến phân. Công bố đầu tiên là của Peruzzo et al. (2014). \\[2pt]
Coupled cluster đơn nhất (UCCSD) & Unitary coupled cluster & Biến thể e\textasciicircum{}(T - T$^\dagger$) của coupled cluster; vì là toán tử đơn nhất nên có thể chuyển thành mạch lượng tử. Đây là ansatz ``mang tính hóa học'' phổ biến nhất cho VQE. \\[2pt]
Ánh xạ Jordan–Wigner / Bravyi–Kitaev & Jordan–Wigner / Bravyi–Kitaev mapping & Quy tắc chuyển toán tử sinh, hủy của fermion thành tích các toán tử Pauli trên qubit mà vẫn giữ đúng hệ thức phản giao hoán. Với Jordan–Wigner, mỗi spin-orbital cần một qubit (Li trong STO-3G cần 10 qubit). \\[2pt]
Ước lượng pha lượng tử (QPE) & Quantum phase estimation & Thuật toán lượng tử tính trị riêng (năng lượng) với độ chính xác cao nếu có trạng thái ban đầu đủ tốt. Nó đòi hỏi máy tính lượng tử chịu lỗi, hiện chưa khả thi ở quy mô của các bài toán hóa học thực tế. \\[2pt]
Lợi thế lượng tử & Quantum advantage & Tình huống máy lượng tử giải một bài toán tốt hơn hẳn máy cổ điển tốt nhất. Với bài toán năng lượng trạng thái cơ bản trong hóa học, bằng chứng về lợi thế ``theo hàm mũ'' vẫn đang được tranh luận (Lee et al., 2023). \\[2pt]
Nhóm tái chuẩn hóa ma trận mật độ (DMRG) & Density matrix renormalization group & Thuật toán do White đề xuất năm 1992, tối ưu hàm sóng dạng MPS bằng cách lặp qua từng orbital và chỉ giữ lại các trạng thái quan trọng nhất. Rất mạnh với hệ có cấu trúc gần ``một chiều'' và với tương quan tĩnh mạnh (nhiều hợp chất kim loại chuyển tiếp). \\[2pt]
Trạng thái tích ma trận (MPS); toán tử tích ma trận (MPO); phân tích giá trị kỳ dị (SVD) & Matrix product state; matrix product operator; singular value decomposition & MPS viết các hệ số của hàm sóng thành tích của một chuỗi ma trận nhỏ, mỗi ma trận ứng với một orbital; MPO là dạng tương tự cho Hamiltonian. SVD là công cụ dùng để cắt bỏ phần ít quan trọng khi nén hàm sóng. \\[2pt]
Số chiều liên kết & Bond dimension (M hoặc D) & Kích thước của các ma trận trong MPS, đóng vai trò ``núm vặn'' độ chính xác của DMRG. Tăng M thì chính xác hơn và tốn hơn; người ta thường ngoại suy theo M để ước lượng giới hạn chính xác. \\[2pt]
Mạng tenxơ & Tensor network & Họ các cách biểu diễn một tenxơ hoặc hàm sóng khổng lồ bằng một mạng các tenxơ nhỏ nối với nhau; MPS là trường hợp đơn giản nhất (dạng chuỗi). Đây là ngôn ngữ chung giữa DMRG, mô phỏng mạch lượng tử và một phần của học máy. \\[2pt]
\bottomrule\end{longtable}
